% sections/privacy.tex — Privacy-Preserving Fingerprint Module
\section{Privacy-Preserving Fingerprint Module}
\label{sec:privacy}

While \system{} can generate reference knowledge graphs from scratch
using configured distributions, a key scenario involves calibrating
the generative model to match real enterprise data.  The fingerprint
module enables this by extracting \emph{distributional summaries} under
formal privacy guarantees, bridging the gap between the forward
generation paradigm (\Cref{sec:gt_forward}) and the statistical
characteristics of production data.

% ─────────────────────────────────────────────────────────────────────
\subsection{Privacy Mechanisms}
\label{sec:priv_mechanisms}

The module combines two complementary techniques:

\paragraph{Differential Privacy ($\varepsilon$-DP).}
Each extracted statistic is protected by calibrated Laplace noise:
\begin{equation}
\tilde{s} = s + \text{Lap}\!\left(\frac{\Delta s}{\varepsilon}\right)
\label{eq:dp_laplace}
\end{equation}
where $s$ is the true statistic, $\Delta s$ is its sensitivity, and
$\varepsilon$ is the privacy budget.  A privacy accountant tracks
cumulative $\varepsilon$ expenditure.

Four privacy levels provide pre-configured $\varepsilon$ values:

\begin{center}
\begin{tabular}{llccc}
\toprule
\textbf{Level} & \textbf{$\varepsilon$} & \textbf{$k$-anon}
  & \textbf{Privacy} & \textbf{Utility} \\
\midrule
Minimal  & 5.0 & 3  & Low    & High \\
Standard & 1.0 & 5  & Medium & Medium \\
High     & 0.5 & 10 & High   & Reduced \\
Maximum  & 0.1 & 20 & Maximum & Low \\
\bottomrule
\end{tabular}
\end{center}

\paragraph{$k$-Anonymity.}
Categorical distributions with fewer than $k$ occurrences are suppressed.

% ─────────────────────────────────────────────────────────────────────
\subsection{Fingerprint Extraction and Synthesis}
\label{sec:priv_extraction}

The \texttt{FingerprintExtractor} processes input data in a streaming
fashion through seven components: schema extraction, statistics
extraction, correlation extraction, integrity extraction, rule
extraction, anomaly extraction, and privacy audit.  Each numeric
statistic is noised via \Cref{eq:dp_laplace}.

The fingerprint is stored as a ZIP archive (\texttt{.dsf}) containing
manifest, schema, statistics, correlations, integrity, rules, anomalies,
and privacy audit components.

The \texttt{ConfigSynthesizer} reads a \texttt{.dsf} fingerprint and
generates a complete \system{} YAML configuration through scale
estimation, distribution mapping, copula reconstruction, anomaly
calibration, and seed assignment.  This two-phase architecture
(extract $\to$ synthesize) cleanly separates the \emph{privacy boundary}
from the \emph{generation process}.

% ─────────────────────────────────────────────────────────────────────
\subsection{Privacy-Utility Tradeoff}
\label{sec:priv_tradeoff}

\begin{proposition}[Utility bound]
\label{prop:utility}
For $n$ records with true mean $\mu$ and column range $R$:
$\E[|\tilde{\mu} - \mu|] = R/(n\varepsilon)$.
For $n = 10{,}000$, $R = 10^6$, $\varepsilon = 1.0$:
$\E[|\tilde{\mu} - \mu|] = 100$, which is $<0.01\%$ of the
range---negligible for distribution fitting.
\end{proposition}

This demonstrates that for enterprise-scale datasets ($n > 10{,}000$),
even standard privacy levels introduce negligible utility loss,
making the fingerprint approach practical for calibrating reference
knowledge graphs against real-world statistical properties while
maintaining formal privacy guarantees.

% ─────────────────────────────────────────────────────────────────────
\subsection{Workflow}
\label{sec:priv_workflow}

\begin{enumerate}
  \item \textbf{Extract:}
    \texttt{datasynth-data fingerprint extract -{}-input data.csv -{}-output fp.dsf -{}-privacy-level standard}
  \item \textbf{Validate:}
    \texttt{datasynth-data fingerprint validate fp.dsf}
  \item \textbf{Generate:}
    \texttt{datasynth-data generate -{}-fingerprint fp.dsf -{}-output ./synthetic/}
  \item \textbf{Evaluate:}
    \texttt{datasynth-data fingerprint evaluate -{}-fingerprint fp.dsf -{}-synthetic ./synthetic/}
\end{enumerate}

The fingerprint file can be safely shared across organizational
boundaries, enabling generation of calibrated reference knowledge
graphs without exposing individual records.
